\chapter{Ruang Lengkap: Baire, Ascoli, Stone--Weierstrass}
\label{ch:b3:complete}

\omterm{def:b3:complete:complete}{Kelengkapan} --- yakni setiap barisan Cauchy konvergen --- adalah
sifat yang membuat analisis dapat \emph{menghasilkan} objek: titik tetap
kontraksi, jumlah deret, penyelesaian persamaan yang diperoleh
sebagai limit. Bab ini merakit ketiga mesin keberadaan
terbesar pada teori metrik. \emph{Teorema Baire} menunjukkan bahwa
ruang \omterm{def:b3:complete:complete}{lengkap} tak mungkin menjadi gabungan terbilang kepingan yang
dapat diabaikan, dan ia memanggil objek (yaitu fungsi \omterm{def:b3:topology:continuity}{kontinu} yang tak
dapat diturunkan di mana pun!) keluar dari penalaran bergaya kardinalitas semata.
\emph{Arzelà--Ascoli} mengenali himpunan bagian \omterm{def:b3:topology:compact}{kompak} dari
$\mathcal C(K)$ dan menjadi kuda beban kekompakan dalam analisis --- dan
soal akhir pekan memakainya untuk membuktikan teorema keberadaan Peano bagi
persamaan diferensial. \emph{Stone--Weierstrass} menunjukkan bahwa
polinomial, dan banyak hal lain, padat di $\mathcal C(K)$: sehingga
penghampiran menjadi pemeriksaan aljabar. Di sepanjang jalan kita
membangun \omterm{thm:b3:complete:completion}{pelengkapan} dan membuktikan teorema perluasan bagi
\omterm{def:b3:topology:continuity}{pemetaan kontinu} seragam, yaitu roti sehari-hari
\cref{ch:b3:lp,ch:b3:hilbert,ch:b3:fouriertransform}.

\section{Ruang lengkap, pelengkapan, perluasan}

\begin{definition}\label{def:b3:complete:complete}
Sebuah ruang metrik disebut \emph{lengkap}\index{ruang metrik lengkap} jika
setiap barisan Cauchy konvergen (Tahun ke-2: $\R^n$ lengkap; dan
$\mathcal C(\intcc01)$ dengan $\norm\cdot_\infty$ lengkap).
Himpunan bagian tertutup dari ruang lengkap bersifat lengkap; sedangkan himpunan bagian lengkap
dari sembarang ruang metrik bersifat tertutup.
\end{definition}

\begin{proof}
Untuk kedua pernyataannya: barisan Cauchy pada $F$ yang tertutup
konvergen di $X$, dan limitnya, yang melekat pada $F$, terletak di $F$; sedangkan
barisan yang konvergen di $X$ pada $A$ yang \omterm{def:b3:complete:complete}{lengkap} bersifat Cauchy, sehingga
konvergen di $A$, dan limitnya tunggal.
\end{proof}

\begin{theorem}[Perluasan pemetaan kontinu seragam]
\label{thm:b3:complete:extension}
Misalkan $D \subseteq X$ padat, $Y$ \omterm{def:b3:complete:complete}{lengkap}, dan $f \colon D \to
Y$ \omterm{def:b3:topology:continuity}{kontinu} seragam. Maka $f$ dapat diperluas secara \emph{tunggal} menjadi
$\bar f \colon X \to Y$ yang \omterm{def:b3:topology:continuity}{kontinu}, dan $\bar f$ \omterm{def:b3:topology:continuity}{kontinu}
seragam.\index{teorema perluasan (kekontinuan seragam)}
\end{theorem}

\begin{proof}
Ketunggalannya: dua perluasan \omterm{def:b3:topology:continuity}{kontinu} bersesuaian pada $D$ yang padat,
sehingga di mana-mana (sebab himpunan kesesuaiannya $\{g = h\}$ tertutup:
yakni prapeta diagonal tertutup di bawah $x\mapsto(g(x), h(x))$).
Keberadaannya: untuk $x \in X$ ambillah $d_n \to x$ dengan $d_n \in D$.
Barisan $(f(d_n))$ bersifat Cauchy: diberikan $\varepsilon$, \omterm{def:b3:topology:continuity}{kekontinuan}
seragamnya menyediakan $\delta$ dengan $d(u,v) < \delta \Rightarrow
d(f(u), f(v)) < \varepsilon$, sedangkan $(d_n)$ bersifat Cauchy. Definisikan
$\bar f(x) = \lim f(d_n)$; limitnya tidak bergantung pada
barisan yang dipilih (selipkan dua di antaranya). Lalu $\bar f$ memperluas $f$
(lewat barisan konstan) dan mewarisi modulus \omterm{def:b3:topology:continuity}{kekontinuannya}: sebab jika
$d(x, x') < \delta$, menghampiri keduanya dengan titik $D$ yang berjarak
$< \frac{\delta - d(x,x')}2$ memberikan $d(\bar f(x), \bar
f(x')) \leq \varepsilon$ pada limitnya --- jadi $\bar f$ \omterm{def:b3:topology:continuity}{kontinu}
seragam.
\end{proof}

\begin{theorem}[Pelengkapan]\label{thm:b3:complete:completion}
Setiap ruang metrik $X$ terbenam secara isometrik sebagai himpunan bagian padat
dari sebuah \omterm{def:b3:complete:complete}{ruang metrik lengkap} $\hat X$, yang tunggal sampai pada isometri yang menetapkan
$X$ titik demi titik: yaitu \emph{pelengkapan}\index{pelengkapan ruang
metrik}nya.
\end{theorem}

\begin{proof}
\emph{Keberadaan.} Misalkan $\mathcal X$ himpunan barisan Cauchy
di $X$, dengan pseudojarak
\[
D\bigl((x_n), (y_n)\bigr) = \lim_n d(x_n, y_n),
\]
dengan limitnya ada karena $\abs{d(x_n, y_n) - d(x_m, y_m)}
\leq d(x_n, x_m) + d(y_n, y_m)$ membuat barisan realnya
Cauchy. Misalkan $\hat X = \mathcal X/{\sim}$, dengan
mengenali barisan yang berjarak-$D$ sama dengan $0$; lalu $D$ turun menjadi
sebuah jarak. Benamkan $X$ lewat barisan konstan: sebuah isometri, dengan
peta yang padat (sebab barisan Cauchy dihampiri secara $D$ oleh
barisan konstan yang dibangun dari sukunya sendiri: $D\bigl((x_n), (x_k)_{\rm
const}\bigr) = \lim_n d(x_n, x_k) \to 0$ ketika $k \to \infty$ menurut
sifat Cauchynya). \omterm{def:b3:complete:complete}{Kelengkapan} $\hat X$: misalkan $(\xi^k)$ Cauchy
di $\hat X$; menurut kepadatannya pilihlah $x_k \in X$ dengan $D(\xi^k, x_k)
\leq 2^{-k}$; maka $(x_k)$ Cauchy di $X$ (lewat ketaksamaan
segitiga melalui $\xi$), sehingga mendefinisikan titik $\xi \in \hat
X$, dan $D(\xi^k, \xi) \leq 2^{-k} + D(x_k, \xi) \to 0$ (sebab
jarak dari konstanta $x_k$ ke kelas $(x_j)_j$ adalah
$\lim_j d(x_k, x_j)$, yang kecil untuk $k$ besar).

\emph{Ketunggalan}: dua \omterm{thm:b3:complete:completion}{pelengkapan} $\hat X_1, \hat X_2$ memuat
$X$ secara padat; identitas $X$, yang merupakan isometri, bersifat \omterm{def:b3:topology:continuity}{kontinu}
seragam, sehingga meluas menjadi $\hat X_1 \to \hat X_2$
(\cref{thm:b3:complete:extension}), yang tetap isometri pada \omterm{def:b3:topology:interior}{himpunan
padat} sehingga di mana-mana; secara simetris pada arah sebaliknya, dan
komposisinya menetapkan $X$ yang padat: jadi keduanya identitas.
\end{proof}

\begin{theorem}[Titik tetap Banach]\label{thm:b3:complete:banach}
Misalkan $X$ \omterm{def:b3:complete:complete}{lengkap} dan tak kosong, serta $f \colon X \to X$ sebuah
\emph{kontraksi}: $d(f(x), f(y)) \leq k\,d(x,y)$ dengan $k < 1$.
Maka $f$ mempunyai titik tetap tunggal $x^*$, dan setiap \omterm{def:b3:groups:action}{orbit}
konvergen kepadanya, dengan laju eksplisit $d(x_n, x^*) \leq
\frac{k^n}{1-k}\,d(x_1, x_0)$.\index{teorema titik tetap Banach}
\end{theorem}

\begin{proof}
(Tahun ke-2 sudah membuktikannya; kita catat ulang argumen dua barisnya demi
kemandirian.) \omterm{def:b3:groups:action}{Orbit} $x_{n+1} = f(x_n)$ memenuhi $d(x_{n+1},
x_n) \leq k^nd(x_1, x_0)$, sehingga bersifat Cauchy (lewat deret geometri);
limitnya $x^*$ bersifat tetap (karena $f$ \omterm{def:b3:topology:continuity}{kontinu}), dan tunggal sebab dua
titik tetap memenuhi $d \leq k\,d$. Lajunya: jumlahkan ekor
geometrinya.
\end{proof}

\begin{example}[Mengusik identitas]
\label{ex:b3:complete:perturbident}
Misalkan $g \colon \R^d \to \R^d$ Lipschitz-$k$ dengan $k < 1$.
Maka $\varphi = \mathrm{id} + g$ merupakan \emph{\omterm{def:b3:topology:continuity}{homeomorfisma}
$\R^d$ pada $\R^d$}. Keinjektifannya, beserta modulus
kuantitatifnya:
\[
\norm{\varphi(x) - \varphi(y)} \geq \norm{x - y} -
\norm{g(x) - g(y)} \geq (1 - k)\norm{x - y} .
\]
Kesurjektifannya adalah teorema titik tetap: sebab menyelesaikan
$\varphi(x) = y$ berarti $x = y - g(x)$, sedangkan $x \mapsto y -
g(x)$ merupakan kontraksi-$k$ pada $\R^d$ yang \omterm{def:b3:complete:complete}{lengkap} --- jadi ada
penyelesaian tunggal $x = \psi(y)$ untuk setiap $y$.
Ketaksamaan yang ditampilkan itu membuat invers $\psi$ bersifat Lipschitz dengan
konstanta $\frac1{1-k}$: jadi sebuah \omterm{def:b3:topology:continuity}{homeomorfisma}, dengan batas
eksplisit pada kedua modulusnya. Pernyataan yang tampak sepele ini adalah
mesin di dalam teorema fungsi invers
(\cref{ch:b3:submanifolds}): sebab di dekat titik tempat $Df$
terbalikkan, $f$ \emph{memang} berupa pemetaan linear terbalikkan ditambah
usikan Lipschitz yang kecil, dan contoh hari ini mengerjakan
sisanya. Ia juga mengukur ketegaran numerik: sistem
yang diusik kurang dari margin inversnya tetap dapat
diselesaikan, dengan penyelesaiannya bergeser paling banyak
$\frac{1}{1-k}$ kali usikannya.
\end{example}

\section{Teorema Baire}

\begin{theorem}[Baire]\label{thm:b3:complete:baire}
Di dalam \omterm{def:b3:complete:complete}{ruang metrik lengkap}, irisan terbilang \omterm{def:b3:topology:topology}{himpunan terbuka} yang
padat tetap padat. Setara dengan itu: jika $X = \bigcup_{n}F_n$ dengan
setiap $F_n$ tertutup, maka ada $F_n$ yang \omterm{def:b3:topology:interior}{interiornya} tak
kosong.\index{teorema Baire}
\end{theorem}

\begin{proof}
Misalkan $(U_n)$ \omterm{def:b3:topology:topology}{himpunan terbuka} yang padat dan $B_0 = B(x_0, r_0)$ sembarang bola
terbuka; kita cari titik $\bigcap U_n$ di dalam $B_0$. Secara induktif:
karena $U_{n}$ padat dan terbuka, ia memotong bola terbuka $B_{n-1}$
pada sebuah \omterm{def:b3:topology:topology}{himpunan terbuka}, yang memuat bola tertutup $\bar B(x_n, r_n)$
dengan $0 < r_n \leq r_{n-1}/2$ dan $\bar B(x_n, r_n) \subseteq
B_{n-1}\cap U_n$. Pusatnya membentuk barisan Cauchy (sebab $x_m \in
B_n$ untuk $m \geq n$, dan jari-jarinya $\to 0$); lalu limitnya $x$ terletak di
setiap $\bar B(x_n, r_n)$ (karena tertutup), sehingga di setiap $U_n$ dan
di $B_0$. Untuk bentuk keduanya: jika tak ada $F_n$ yang berinterior, maka
$U_n = X \setminus F_n$ terbuka dan padat, dan sebuah titik
$\bigcap U_n$ lolos dari $\bigcup F_n = X$: mustahil.
\end{proof}

\begin{remark}\label{rem:b3:complete:meagre}
Kosakatanya: sebuah himpunan disebut \emph{tak padat di mana pun} jika \omterm{def:b3:topology:interior}{penutupnya}
berinterior kosong, dan disebut \emph{kurus}\index{himpunan kurus} (kategori pertama)
jika ia gabungan terbilang himpunan yang tak padat di mana pun. Baire mengatakan:
\emph{\omterm{def:b3:complete:complete}{ruang metrik lengkap} tidak \omterm{rem:b3:complete:meagre}{kurus} di dalam dirinya sendiri}, dan
komplemen \omterm{rem:b3:complete:meagre}{himpunan kurus} bersifat padat. ``\omterm{rem:b3:complete:meagre}{Kurus}'' adalah gagasan
kekecilan yang ortogonal terhadap ukuran (\cref{ch:b3:measure} akan
menghasilkan \omterm{rem:b3:complete:meagre}{himpunan kurus} yang berukuran penuh), dan argumen Baire membuktikan
\emph{keberadaan lewat kelimpahan}: untuk menampilkan satu objek tanpa
sifat P, tunjukkan bahwa objek berciri P membentuk himpunan yang \omterm{rem:b3:complete:meagre}{kurus}.
\end{remark}

\begin{corollary}\label{cor:b3:complete:bairecors}
(a) $\R$ tak terbilang. (b) $\Q$ bukan irisan terbilang
himpunan bagian terbuka $\R$, dan \omterm{def:b3:complete:complete}{ruang metrik lengkap} tak kosong
tanpa titik terisolasi bersifat tak terbilang.
\end{corollary}

\begin{proof}
(a) Jika $\R = \bigcup_{x}\{x\}$ atas himpunan terbilang, maka suatu
himpunan tunggal akan berinterior. (b) Jika $\Q = \bigcap_n V_n$ dengan $V_n$
terbuka (dan tentu padat, sebab $\supseteq \Q$), maka himpunan
$V_n$ beserta komplemen $\R\setminus\{q\}$ dengan $q \in \Q$ membentuk
keluarga terbilang \omterm{def:b3:topology:topology}{himpunan terbuka} padat yang irisannya kosong ---
dan itu bertentangan dengan Baire. Jika $X$ \omterm{def:b3:complete:complete}{lengkap} tanpa titik terisolasi
dan terbilang, maka $X = \bigcup_{x \in X}\{x\}$ menampilkannya sebagai
gabungan terbilang himpunan tertutup berinterior kosong (karena tanpa titik
terisolasi): sekali lagi Baire.
\end{proof}

\begin{theorem}[Monster Weierstrass memang ada]
\label{thm:b3:complete:nowherediff}
Ada fungsi \omterm{def:b3:topology:continuity}{kontinu} pada $\intcc01$ yang tak dapat diturunkan di
satu titik pun. Bahkan, himpunan $f \in \mathcal C(\intcc01)$
yang mempunyai turunan (berhingga) walau di satu titik saja bersifat \omterm{rem:b3:complete:meagre}{kurus} di
$(\mathcal C(\intcc01), \norm\cdot_\infty)$.\index{fungsi tak dapat
diturunkan di mana pun}
\end{theorem}

\begin{proof}
Untuk $n \geq 1$ misalkan
\[
F_n = \Bigl\{f : \exists x \in \intcc01,\
\forall h \neq 0 \text{ dengan } x + h \in \intcc01,\
\abs{f(x+h) - f(x)} \leq n\abs h \Bigr\}.
\]
Jika $f$ dapat diturunkan di $x$, maka $f \in F_n$ untuk suatu $n$:
sebab hasil bagi $\abs{f(x+h)-f(x)}/\abs h$ terbatas untuk $\abs h
\leq \delta$ (menurut keterdiferensialannya: ia menuju
$\abs{f'(x)}$) dan terbatas oleh $2\norm f_\infty/\delta$ untuk
$\abs h \geq \delta$. Jadi $\bigcup F_n$ memuat
semua fungsi yang dapat diturunkan di suatu tempat, sehingga cukup ditunjukkan bahwa
setiap $F_n$ tertutup dan berinterior kosong.

\emph{Tertutup}: misalkan $f_k \to f$ secara seragam, dengan $f_k \in F_n$ beserta
saksinya $x_k \to x$ (lewat kekompakan, setelah ekstraksi). Untuk $h$
dengan $x + h \in \intcc01$: pilihlah $h_k \to h$ dengan $x_k + h_k
\in \intcc01$ (misalnya $h_k = h + x - x_k$ yang dipangkas); maka
$\abs{f(x + h) - f(x)} = \lim \abs{f_k(x_k + h_k) - f_k(x_k)}
\leq \lim n\abs{h_k} = n\abs h$, dengan memakai kekonvergenan seragam dan
\omterm{def:b3:topology:continuity}{kekontinuan} $f$ di titik yang bersangkutan: jadi $f \in F_n$.

\emph{\omterm{def:b3:topology:interior}{Interior} kosong}: diberikan $f \in F_n$ dan $\varepsilon > 0$,
kita cari $g$ dengan $\norm{g - f}_\infty \leq \varepsilon$ dan $g
\notin F_n$. Pertama hampiri $f$ dalam $\varepsilon/2$ oleh
fungsi afin sepotong-sepotong $\varphi$ (lewat \omterm{def:b3:topology:continuity}{kekontinuan} seragamnya: interpolasikan pada
kisi yang halus), yang kemiringannya terbatas oleh suatu $M$. Lalu tambahkan gigi gergaji kecil:
$g = \varphi + \frac\varepsilon2\,s_N$, dengan $s_N(x)$ berupa
zigzag berperiode $\frac1N$, beramplitudo $1$ dan berkemiringan $\pm 2N$.
Di setiap $x$, pada salah satu sisinya ada $h$ sekecil apa pun dengan
gigi gergajinya menyumbang kemiringan $\pm 2N\cdot\frac\varepsilon2 =
\pm\varepsilon N$ sepanjang $[x, x+h]$: sehingga hasil bagi selisih
$g$ melampaui $\varepsilon N - M > n$ untuk $N$ besar. Jadi $g \notin
F_n$, pada jarak seragam $\leq \varepsilon$ berapa pun dari $f$.
Kesimpulannya: $\bigcup F_n$ bersifat \omterm{rem:b3:complete:meagre}{kurus}; dan menurut Baire komplemennya
--- yang terdiri atas \omterm{thm:b3:complete:nowherediff}{fungsi tak dapat diturunkan di mana pun} --- bersifat padat di
$\mathcal C(\intcc01)$: jadi fungsi semacam itu ada \emph{dengan
berlimpah}.
\end{proof}

\section{Arzelà--Ascoli}

Sepanjang bagian ini, $K$ adalah ruang metrik \emph{\omterm{def:b3:topology:compact}{kompak}} dan $\mathcal
C(K) = \mathcal C(K, \R^d)$, dengan $\norm f_\infty = \sup_K
\norm{f(x)}$: sebuah ruang \omterm{def:b3:complete:complete}{lengkap} (sebab limit seragam fungsi \omterm{def:b3:topology:continuity}{kontinu}
tetap \omterm{def:b3:topology:continuity}{kontinu} --- Tahun ke-2).

\begin{definition}\label{def:b3:complete:equicontinuous}
Sebuah keluarga $\mathcal F \subseteq \mathcal C(K)$ disebut
\emph{ekuikontinu}\index{keekuikontinuan} jika untuk setiap
$\varepsilon > 0$ ada $\delta > 0$ sedemikian sehingga
\[
d(x, y) < \delta \implies \norm{f(x) - f(y)} < \varepsilon
\quad\text{untuk \emph{setiap} } f \in \mathcal F
\]
(yakni satu $\delta$ untuk seluruh keluarganya --- misalnya keluarga apa pun dengan
konstanta Lipschitz bersama, atau modulus Hölder bersama), dan disebut
\emph{terbatas titik demi titik} jika $\sup_{f}\norm{f(x)} < \infty$ untuk
setiap $x$.
\end{definition}

\begin{theorem}[Arzelà--Ascoli]\label{thm:b3:complete:ascoli}
Sebuah himpunan bagian $\mathcal F \subseteq \mathcal C(K)$ bersifat
\omterm{def:b3:topology:compact}{kompak} relatif (yakni \omterm{def:b3:topology:interior}{penutupnya} \omterm{def:b3:topology:compact}{kompak}) jika dan hanya jika ia \omterm{def:b3:complete:equicontinuous}{ekuikontinu} dan
terbatas titik demi titik. Khususnya, setiap \emph{barisan} yang \omterm{def:b3:complete:equicontinuous}{ekuikontinu} dan
terbatas titik demi titik mempunyai subbarisan yang konvergen
seragam.\index{teorema Arzelà--Ascoli}
\end{theorem}

\begin{proof}
($\Leftarrow$) Misalkan $(f_k)$ sebuah barisan di $\mathcal F$. Ruang $K$
metrik \omterm{def:b3:topology:compact}{kompak} bersifat \emph{\omterm{def:b3:galois:separable}{separabel}}: sebab untuk setiap $n$, berhingga banyak
bola berjari-jari $\frac1n$ menyelimuti $K$
(\cref{thm:b3:topology:metriccompact}); dan pusatnya membentuk
\omterm{def:b3:topology:interior}{himpunan padat} terbilang $D = \{x_1, x_2, \dots\}$. Menurut keterbatasan titik demi
titik dan Bolzano--Weierstrass, ekstraklah berturut-turut
subbarisan yang konvergen di $x_1$, lalu juga di $x_2$, dan seterusnya, lalu
ambil subbarisan \emph{diagonalnya} $(g_j)$: ia konvergen di
setiap titik $D$. \omterm{def:b3:complete:equicontinuous}{Keekuikontinuannya} meningkatkan hal ini menjadi Cauchy
seragam: diberikan $\varepsilon$, ambillah $\delta$ seperti pada
definisinya, selimuti $K$ dengan berhingga bola $B(x_i, \delta)$
dengan $x_i \in D$ ($i \leq m$), lalu pilihlah $J$ sedemikian besar sehingga
$\norm{g_j(x_i) - g_{j'}(x_i)} < \varepsilon$ untuk $j, j' \geq
J$ dan $i \leq m$. Untuk sembarang $x \in B(x_i, \delta)$:
\[
\norm{g_j(x) - g_{j'}(x)}
\leq \norm{g_j(x) - g_j(x_i)} + \norm{g_j(x_i) - g_{j'}(x_i)}
+ \norm{g_{j'}(x_i) - g_{j'}(x)} < 3\varepsilon .
\]
Jadi $(g_j)$ Cauchy seragam, sehingga konvergen di
$\mathcal C(K)$ yang \omterm{def:b3:complete:complete}{lengkap}. Dengan demikian setiap barisan di $\mathcal F$ mempunyai
subbarisan konvergen: jadi $\bar{\mathcal F}$ \omterm{def:b3:topology:compact}{kompak} (secara barisan,
sehingga menurut \cref{thm:b3:topology:metriccompact} \omterm{def:b3:topology:compact}{kompak}).

($\Rightarrow$) Jika $\bar{\mathcal F}$ \omterm{def:b3:topology:compact}{kompak}: maka keterbatasan titik demi
titiknya jelas (sebab pemasukan nilai bersifat \omterm{def:b3:topology:continuity}{kontinu}). Untuk
\omterm{def:b3:complete:equicontinuous}{keekuikontinuannya}, selimuti $\mathcal F$ dengan berhingga bola
$B(f_i, \varepsilon)$ di $\mathcal C(K)$; setiap $f_i$ bersifat
\omterm{def:b3:topology:continuity}{kontinu} seragam (menurut Heine,
\cref{cor:b3:topology:heineborel}), sehingga memberikan $\delta$ bersama
untuk $i \leq m$; lalu untuk $f \in B(f_i, \varepsilon)$ dan $d(x,y)
< \delta$: $\norm{f(x)-f(y)} \leq 2\varepsilon +
\norm{f_i(x)-f_i(y)} < 3\varepsilon$.
\end{proof}

\begin{example}\label{ex:b3:complete:ascoliexamples}
Bola satuan tertutup di $\mathcal C(\intcc01)$ \emph{tidak}
\omterm{def:b3:topology:compact}{kompak} (sebab $f_n(x) = x^n$ tak punya subbarisan yang konvergen
seragam: limit titik demi titiknya tak \omterm{def:b3:topology:continuity}{kontinu}), dan memang
$(x^n)$ tidak \omterm{def:b3:complete:equicontinuous}{ekuikontinu} di $1$. Sebaliknya $\{f : \norm
f_\infty \leq 1,\ \operatorname{Lip}(f) \leq 1\}$ bersifat \omterm{def:b3:topology:compact}{kompak}:
sebab terbatas dan \omterm{def:b3:complete:equicontinuous}{ekuikontinu} lewat Lipschitz-$1$, serta tertutup. Ascoli
menjelaskan \emph{mengapa} kekompakan gagal pada dimensi tak berhingga
(Riesz, Tahun ke-2) dan apa yang harus ditambahkan untuk memulihkannya: yaitu
modulus \omterm{def:b3:topology:continuity}{kekontinuan} yang seragam.
\end{example}

\section{Stone--Weierstrass}

\begin{lemma}[Dini]\label{lem:b3:complete:dini}
Misalkan $K$ \omterm{def:b3:topology:compact}{kompak} dan $(f_n)$ barisan \emph{monoton} berisi
fungsi real \omterm{def:b3:topology:continuity}{kontinu} yang konvergen \emph{titik demi titik} ke
$f$ yang \omterm{def:b3:topology:continuity}{kontinu}. Maka kekonvergenannya bersifat
seragam.\index{teorema Dini}
\end{lemma}

\begin{proof}
Katakanlah $f_n \uparrow f$; misalkan $g_n = f - f_n \downarrow 0$,
yang \omterm{def:b3:topology:continuity}{kontinu}. Diberikan $\varepsilon$, \omterm{def:b3:topology:topology}{himpunan terbuka} $U_n = \{g_n <
\varepsilon\}$ naik dan menyelimuti $K$ (menurut kekonvergenan titik demi titiknya);
ekstraklah subselimut berhingga: maka $K = U_{n_0}$ untuk suatu $n_0$
(karena keluarganya naik), yakni $0 \leq g_n < \varepsilon$
di mana-mana untuk $n \geq n_0$.
\end{proof}

\begin{lemma}\label{lem:b3:complete:sqrt}
Ada barisan \emph{polinomial} $u_n$ dengan $u_n(t) \to
\sqrt t$ secara seragam pada $\intcc01$.
\end{lemma}

\begin{proof}
Definisikan $u_0 = 0$ dan $u_{n+1}(t) = u_n(t) + \frac12\bigl(t -
u_n(t)^2\bigr)$: semuanya polinomial. Dengan induksi, $0 \leq u_n(t) \leq
\sqrt t$ pada $\intcc01$: sebab dengan mengandaikannya untuk $n$,
\[
\sqrt t - u_{n+1}(t) = \bigl(\sqrt t - u_n(t)\bigr)
\Bigl(1 - \tfrac12\bigl(\sqrt t + u_n(t)\bigr)\Bigr) \geq 0,
\]
karena $\sqrt t + u_n \leq 2$; lalu $u_{n+1} \geq u_n \geq 0$.
Jadi $(u_n(t))$ tak turun dan terbatas oleh $\sqrt t$: sehingga ia
konvergen titik demi titik, dan limitnya $\ell(t)$ memenuhi $\ell =
\ell + \frac12(t - \ell^2)$: jadi $\ell(t) = \sqrt t$, yang \omterm{def:b3:topology:continuity}{kontinu}.
Lalu Dini (\cref{lem:b3:complete:dini}) meningkatkannya menjadi seragam.
\end{proof}

\begin{theorem}[Stone--Weierstrass, versi real]
\label{thm:b3:complete:stoneweierstrass}
Misalkan $K$ sebuah \omterm{def:b3:topology:compact}{ruang kompak} (\omterm{def:b3:topology:hausdorff}{Hausdorff}) dan $\mathcal A \subseteq
\mathcal C(K, \R)$ sebuah \emph{subaljabar} (yang stabil terhadap penjumlahan,
perkalian, dan kelipatan skalar) yang \emph{memuat konstanta}
serta \emph{memisahkan titik} (yakni untuk $x \neq y$ ada $f \in
\mathcal A$ dengan $f(x) \neq f(y)$). Maka $\mathcal A$ padat di
$(\mathcal C(K,\R), \norm\cdot_\infty)$.\index{teorema
Stone--Weierstrass}
\end{theorem}

\begin{proof}
Misalkan $\bar{\mathcal A}$ \omterm{def:b3:topology:interior}{penutupnya}, yang kembali menjadi aljabar
(sebab hasil kali limit seragam pada himpunan terbatas tetap konvergen).

\emph{Langkah 1: $\bar{\mathcal A}$ merupakan latis}, yakni stabil
terhadap $\max$ dan $\min$. Karena $\max(f,g) = \frac{f + g +
\abs{f-g}}2$ dan demikian pula $\min$, cukup dibuktikan bahwa $f \in
\bar{\mathcal A} \Rightarrow \abs f \in \bar{\mathcal A}$: dengan
$M = \norm f_\infty > 0$ berlaku $\abs f = M\sqrt{(f/M)^2}$, dan
\cref{lem:b3:complete:sqrt} memberikan polinomial $u_n$ dengan
$u_n\bigl((f/M)^2\bigr) \to \abs f/M$ secara seragam; sedangkan polinomial atas
anggota aljabarnya (dengan suku konstan: konstantanya memang
ada) tetap berada di $\bar{\mathcal A}$.

\emph{Langkah 2: interpolasi dua titik.} Untuk $x \neq y$ dan $a,
b \in \R$, ada $g \in \mathcal A$ dengan $g(x) = a$, $g(y) = b$:
ambil $h$ yang memisahkan $x, y$ lalu tetapkan $g = a + (b -
a)\frac{h - h(x)}{h(y) - h(x)}$.

\emph{Langkah 3.} Misalkan $f \in \mathcal C(K)$ dan $\varepsilon > 0$.
Untuk setiap pasangan $x, y$ pilihlah $g_{x,y} \in \mathcal A$ dengan
$g_{x,y}(x) = f(x)$, $g_{x,y}(y) = f(y)$ (lewat Langkah 2; sedangkan untuk $x = y$
ambillah fungsi konstan $g_{x,x} = f(x)$).
Tetapkan $x$: untuk setiap $y$, \omterm{def:b3:topology:topology}{himpunan terbuka} $V_y = \{g_{x,y} < f +
\varepsilon\}$ memuat $y$; lalu kekompakannya mengekstrak $y_1, \dots,
y_m$ dengan $K = \bigcup V_{y_j}$, dan $h_x = \min_j g_{x, y_j}
\in \bar{\mathcal A}$ (Langkah 1) memenuhi $h_x < f +
\varepsilon$ di mana-mana serta $h_x(x) = f(x)$. Kini ragamkan $x$: $W_x =
\{h_x > f - \varepsilon\}$ terbuka dan memuat $x$; ekstraklah $x_1,
\dots, x_l$ yang menyelimuti $K$, lalu $h = \max_i h_{x_i} \in
\bar{\mathcal A}$ memenuhi $f - \varepsilon < h < f +
\varepsilon$: sehingga $\norm{f - h}_\infty \leq \varepsilon$. Karena itu
$f \in \bar{\mathcal A}$.
\end{proof}

\begin{corollary}\label{cor:b3:complete:weierstrass}
(a) (\emph{Weierstrass}) Polinomial padat di $\mathcal
C(\intcc ab, \R)$; dan polinomial dalam $n$ variabel padat di
$\mathcal C(K, \R)$ untuk $K \subseteq \R^n$ yang \omterm{def:b3:topology:compact}{kompak}.
(b) (\emph{Versi kompleks}) Jika $\mathcal A \subseteq \mathcal
C(K, \C)$ sebuah subaljabar yang memuat konstanta, memisahkan
titik, dan \emph{stabil terhadap konjugasi}, maka ia padat.
(c) (\emph{Versi trigonometri}) Polinomial trigonometri
$\sum_{\abs n \leq N}c_n\eu^{\iu n t}$ padat di ruang berisi
fungsi \omterm{def:b3:topology:continuity}{kontinu} berperiode $2\pi$ dengan
$\norm\cdot_\infty$.\index{teorema penghampiran Weierstrass}
\end{corollary}

\begin{proof}
(a) Polinomial membentuk aljabar yang memuat konstanta; dan fungsi
koordinatnya memisahkan titik $\R^n$.
(b) Bagian real dan imajiner $\frac{f + \bar f}2$ dan
$\frac{f - \bar f}{2\iu}$ dari anggota $\mathcal A$ membentuk
aljabar real $\mathcal A_\R \subseteq \mathcal C(K, \R)$ yang memuat
konstanta; ia memisahkan titik (sebab $f(x) \ne f(y)$ memaksa
$\operatorname{Re}f$ atau $\operatorname{Im}f$ memisahkannya).
Lalu terapkan teorema realnya dan gabungkan kembali.
(c) Pandanglah fungsi berperiode $2\pi$ sebagai $\mathcal C(S^1, \C)$
(dengan $S^1 = \R/2\pi\Z$ yang \omterm{def:b3:topology:compact}{kompak}:
\cref{exo:b3:topology:5}); maka aljabar yang dibangun oleh
$\eu^{\iu t}, \eu^{-\iu t}$ dan konstanta bersifat stabil terhadap
konjugasi dan memisahkan titik lingkarannya (sebab $\eu^{\iu t}$
injektif di sana). Lalu terapkan (b).
\end{proof}

\begin{remark}\label{rem:b3:complete:swremark}
Versi trigonometrinya memperbaiki, dan menyamaratakan secara luas,
celah yang ditinggalkan bab Fourier pada Tahun ke-2: sebab kepadatan polinomial
trigonometri di $(\mathcal C(S^1), \norm\cdot_2)$ menyusul secara
a fortiori (karena $\norm\cdot_2 \leq \norm\cdot_\infty$ sampai pada
konstanta penormalannya), dan itulah yang akan menjadikan sistem Fourier sebuah
\emph{basis} ortonormal pada \cref{ch:b3:hilbert}, sehingga akhirnya
membuktikan Parseval dalam keumuman penuh.
\end{remark}

\section{Latihan}

\begin{exercise}[$\star$]\label{exo:b3:complete:1}
(a) Tunjukkan bahwa $\mathcal C(\intcc01, \R)$ dengan
$\norm f_1 = \int_0^1\abs f$ \emph{tidak} \omterm{def:b3:complete:complete}{lengkap}: sebab
fungsi $f_n$ yang bernilai $0$ pada $\intcc0{\frac12}$, $1$ pada
$\intcc{\frac12 + \frac1n}1$, dan afin di antaranya, bersifat
Cauchy-$\norm\cdot_1$ tanpa limit yang \omterm{def:b3:topology:continuity}{kontinu}.
(b) Tunjukkan bahwa ruang bernorma yang setiap deret konvergen
mutlaknya konvergen bersifat \omterm{def:b3:complete:complete}{lengkap}. \emph{(Ekstraklah dari sebuah
barisan Cauchy sebuah subbarisan dengan $\norm{x_{n_{k+1}} -
x_{n_k}} \leq 2^{-k}$.)}
\end{exercise}

\begin{exercise}[$\star$]\label{exo:b3:complete:2}
Dengan memakai Baire: (a) tunjukkan bahwa ruang bernorma yang \omterm{def:b3:complete:complete}{lengkap} tak punya
basis (aljabar) yang terbilang --- lalu simpulkan bahwa ruang
polinomial tidak \omterm{def:b3:complete:complete}{lengkap} terhadap norma \emph{apa pun}; (b) tunjukkan bahwa jika
barisan fungsi \omterm{def:b3:topology:continuity}{kontinu} $f_n \colon \R \to \R$ konvergen
titik demi titik ke $f$, maka himpunan titik \omterm{def:b3:topology:continuity}{kekontinuan} $f$ bersifat padat.
\emph{(Untuk (b): terimalah atau buktikan bahwa $\Omega_\delta = \{x :
\operatorname{osc}_x f < \delta\}$ terbuka, lalu tunjukkan bahwa ia
padat dengan memakai $F_{N} = \{x: \abs{f_n(x) - f_m(x)} \leq \delta/3\
\forall n,m \geq N\}$, yaitu himpunan tertutup yang menyelimuti $\R$; bekerjalah di
sembarang bola tertutup agar Baire dapat diterapkan.)}
\end{exercise}

\begin{exercise}[$\star\star$]\label{exo:b3:complete:3}
(a) Tunjukkan bahwa $f(x) = \frac12\bigl(x + \frac ax\bigr)$ (dengan $a >
1$) merupakan kontraksi pada $[\sqrt a, +\infty)$ lalu kenali titik
tetapnya --- yakni metode Heron. Taksirlah banyaknya
iterasi untuk ketelitian $10^{-12}$ mulai dari $x_0 = a$,
untuk $a = 2$.
(b) (Persamaan Kepler) Untuk $0 \leq e < 1$ dan $m \in \R$, tunjukkan
bahwa $x = m + e\sin x$ mempunyai penyelesaian tunggal yang bergantung
\omterm{def:b3:topology:continuity}{kontinu} pada $m$.
\end{exercise}

\begin{exercise}[$\star\star$]\label{exo:b3:complete:4}
(a) Misalkan $X$ \omterm{def:b3:complete:complete}{lengkap} dan $f \colon X \to X$ sedemikian sehingga suatu
iterasinya $f^p$ merupakan kontraksi. Tunjukkan bahwa $f$ mempunyai titik
tetap tunggal. Terapannya: operator integral $T$ pada
$\mathcal C(\intcc0a)$ dengan $Tf(x) = \int_0^x f$ memenuhi
$\norm{T^p} \leq a^p/p!$ --- lalu simpulkan bahwa $u = g + \lambda Tu$
\omterm{def:b3:groups:derived}{terselesaikan} untuk setiap $\lambda$.
(b) (Edelstein) Misalkan $K$ \emph{\omterm{def:b3:topology:compact}{kompak}} dan $f\colon K \to K$
dengan $d(f(x), f(y)) < d(x, y)$ untuk $x \neq y$. Tunjukkan bahwa $f$
mempunyai titik tetap tunggal, tetapi laju kontraksinya dapat
hilang: pada $X = [1, +\infty)$ (yang \omterm{def:b3:complete:complete}{lengkap} tetapi tak \omterm{def:b3:topology:compact}{kompak}), $f(x) = x
+ \frac1x$ tak punya titik tetap meskipun jaraknya menurun
sejati.
\end{exercise}

\begin{exercise}[$\star\star$]\label{exo:b3:complete:5}
(a) Dua \omterm{def:b3:topology:continuity}{pemetaan kontinu} ke \omterm{def:b3:topology:hausdorff}{ruang Hausdorff} yang bersesuaian pada
himpunan bagian padat akan bersesuaian di mana-mana; di manakah hal ini dipakai pada
bab ini?
(b) Misalkan $D \subseteq X$ padat dan $f \colon D \to Y$ sebuah
bijeksi isometrik pada himpunan bagian padat dari $Y$ yang \omterm{def:b3:complete:complete}{lengkap}, dengan
$X$ \omterm{def:b3:complete:complete}{lengkap}. Tunjukkan bahwa $f$ meluas menjadi bijeksi isometrik
$X \to Y$. Lalu simpulkan kembali ketunggalan \omterm{thm:b3:complete:completion}{pelengkapannya}.
\end{exercise}

\begin{exercise}[$\star\star$]\label{exo:b3:complete:6}
Manakah di antara keluarga berikut yang \omterm{def:b3:complete:equicontinuous}{ekuikontinu}, terbatas titik demi
titik, dan \omterm{def:b3:topology:compact}{kompak} relatif di $\mathcal C(\intcc01)$?
\[
\{x \mapsto \sin(nx)\}_{n},\qquad
\{x \mapsto x^n\}_n,
\]
\[
\Bigl\{f \in \mathcal C^1 : \norm f_\infty \leq 1,\
\norm{f'}_\infty \leq 1\Bigr\},\qquad
\{f : \operatorname{Lip}(f)\leq 1,\ f(0) = 0\}.
\]
Berilah alasan bagi setiap jawabannya dengan Ascoli atau dengan barisan penyangkal.
\end{exercise}

\begin{exercise}[$\star\star\star$]\label{exo:b3:complete:7}
(Operator integral yang \omterm{def:b3:topology:compact}{kompak}) Misalkan $k \in \mathcal C(\intcc01^2)$
dan, untuk $f \in \mathcal C(\intcc01)$, $Tf(x) =
\int_0^1k(x,y)f(y)\,\dd y$.
(a) Tunjukkan bahwa $T$ memetakan bola satuan $\mathcal C(\intcc01)$
ke himpunan yang \omterm{def:b3:complete:equicontinuous}{ekuikontinu} dan terbatas seragam; lalu simpulkan bahwa $T$
merupakan \emph{operator \omterm{def:b3:topology:compact}{kompak}}: yakni peta himpunan terbatas bersifat
\omterm{def:b3:topology:compact}{kompak} relatif.
(b) Simpulkan bahwa jika $(f_n)$ terbatas, maka $(Tf_n)$ mempunyai subbarisan
yang konvergen seragam, dan bahwa $T$ tak mungkin menjadi bijeksi dengan
invers yang \omterm{def:b3:topology:continuity}{kontinu}. \emph{(Sebab peta bola satuannya akan menjadi
\omterm{def:b3:topology:topology}{persekitaran} \omterm{def:b3:topology:compact}{kompak} $0$ di $\mathcal C(\intcc01)$:
dan itu dilarang teorema Riesz dari Tahun ke-2.)}
\end{exercise}

\begin{exercise}[$\star\star$]\label{exo:b3:complete:8}
Buktikan atau sanggahlah, untuk $f_n \colon \intcc01 \to \R$ yang \omterm{def:b3:topology:continuity}{kontinu}:
(a) $f_n \downarrow 0$ titik demi titik $\Rightarrow$ secara seragam (yakni Dini
--- buktikan ulang); (b) sama tetapi tanpa kemonotonan; (c) sama dengan
kemonotonan tetapi $f$ tak \omterm{def:b3:topology:continuity}{kontinu}; (d) sama dengan kemonotonan,
dengan limit \omterm{def:b3:topology:continuity}{kontinu}, tetapi pada $\intoo01$.
\end{exercise}

\begin{exercise}[$\star\star$]\label{exo:b3:complete:9}
(a) (Momen menentukan) Misalkan $f \in \mathcal C(\intcc01)$ dengan
$\int_0^1 f(x)\,x^n\,\dd x = 0$ untuk setiap $n \in \N$. Tunjukkan $f =
0$. \emph{(Hampiri $f$ secara seragam dengan polinomial lalu hitunglah
$\int f^2$.)}
(b) Tunjukkan bahwa polinomial genap padat di $\mathcal
C(\intcc01)$ tetapi \emph{tidak} di $\mathcal C(\intcc{-1}1)$; di manakah
hipotesis Stone--Weierstrass gagal?
(c) Apakah aljabar yang dibangun oleh $x \mapsto \eu^{\iu x}$ saja
(tanpa $\eu^{-\iu x}$) padat di $\mathcal C(S^1, \C)$?
\emph{(Tinjaulah $\int_0^{2\pi}f(t)\,\eu^{\iu t}\dd t$.)}
\end{exercise}

\begin{exercise}[$\star\star\star$]\label{exo:b3:complete:10}
(Keterbatasan seragam, versi metrik) Misalkan $X$ sebuah ruang metrik
yang \omterm{def:b3:complete:complete}{lengkap} dan $\mathcal F \subseteq \mathcal C(X, \R)$ sebuah
keluarga yang terbatas \emph{titik demi titik}: $\sup_{f \in \mathcal
F}\abs{f(x)} < \infty$ untuk setiap $x$. Tunjukkan bahwa ada
$U \subseteq X$ terbuka dan tak kosong yang di atasnya $\mathcal F$ terbatas
\emph{seragam}: $\sup_{f}\sup_U \abs f < \infty$.
\emph{(Tinjaulah $F_n = \{x : \abs{f(x)} \leq n\ \forall f\}$.)}
Inilah mesin di balik Banach--Steinhaus pada
\cref{ch:b3:banach}.
\end{exercise}

\begin{exercise}[$\star\star$]\label{exo:b3:complete:11}
($\mathcal C^1$ memerlukan normanya sendiri) Pada $E = \mathcal
C^1(\intcc01, \R)$ tinjaulah $\norm f_{\mathcal C^1} = \norm
f_\infty + \norm{f'}_\infty$.
(a) Tunjukkan bahwa $(E, \norm\cdot_{\mathcal C^1})$ \omterm{def:b3:complete:complete}{lengkap}
\emph{(barisan Cauchy-$\mathcal C^1$ mempunyai $f_n \to f$ dan
$f_n' \to g$ secara seragam; kenalilah $g = f'$ dengan mengambil limit
pada $f_n(x) = f_n(0) + \int_0^xf_n'$)}.
(b) Tunjukkan bahwa $(E, \norm\cdot_\infty)$ \emph{tidak}
\omterm{def:b3:complete:complete}{lengkap}: tampilkan limit seragam fungsi $\mathcal C^1$
yang tak dapat diturunkan (misalnya penghampiran mulus dari
$\abs{x - \tfrac12}$).
(c) Simpulkan dari (a), (b) dan lingkaran gagasan pemetaan terbuka
--- atau secara langsung --- bahwa tak ada konstanta $C$ yang memenuhi
$\norm{f'}_\infty \leq C\,\norm f_\infty$ pada $E$: tampilkan
barisan yang menjadi saksinya. Jadi pendiferensialan tak terbatas; dan inilah
jurang di balik \cref{thm:b3:complete:nowherediff}.
\end{exercise}

\begin{exercise}[$\star\star\star$]\label{exo:b3:complete:12}
(Lema Croft) Misalkan $f\colon\intoo0\infty\to\R$ \omterm{def:b3:topology:continuity}{kontinu}
dan andaikan bahwa untuk setiap $x > 0$ berlaku $f(nx) \to 0$ ketika
bilangan bulat $n \to \infty$. Tunjukkan bahwa $f(t) \to 0$ ketika $t \to
+\infty$. \emph{(Tetapkan $\varepsilon > 0$; maka himpunan
\[
F_N = \{x > 0 : \abs{f(nx)} \leq \varepsilon\ \ \forall n
\geq N\}
\]
bersifat tertutup dan menyelimuti $\intoo0\infty$; lalu Baire pada suatu selang
$\intcc ab$ memberikan $N$ dan subselang $\intcc{a'}{b'}
\subseteq F_N$; kemudian dilatasi $\bigl[na', nb'\bigr]$ dengan $n
\geq N$ menyelimuti seluruh \omterm{def:b3:topology:topology}{persekitaran} $+\infty$ begitu
$n(b' - a') \geq a'$.)} Di manakah hipotesis ``untuk setiap
$x$'' (bukan hanya $x$ rasional) dipakai?
\end{exercise}

\section{Soal: teorema keberadaan Peano}

\begin{problem}[{Soal akhir pekan --- keberadaan penyelesaian
$x' = f(t, x)$ tanpa Lipschitz}]\label{pb:b3:complete:1}
Cauchy--Lipschitz (Tahun ke-2; dibuktikan ulang pada \cref{ch:b3:ode})
menuntut $f$ bersifat Lipschitz terhadap $x$. Peano (1890): \emph{\omterm{def:b3:topology:continuity}{kekontinuan}
$f$ sudah melahirkan keberadaannya} --- meski bukan ketunggalannya. Kita
membuktikannya dengan poligon Euler dan Ascoli. Latarnya: $f \colon R
\to \R^d$ \omterm{def:b3:topology:continuity}{kontinu} pada persegi panjang $R = \intcc{t_0 - a}{t_0
+ a}\times \bar B(x_0, b) \subseteq \R\times\R^d$, dengan $M =
\sup_R\norm f$, dan
\[
T = \min\bigl(a,\ b/M\bigr)
\qquad (\text{jika } M = 0 \text{ soalnya trivial}).
\]

\medskip
\noindent\textbf{Bagian I --- Poligon Euler.} Untuk $n \geq 1$,
bagilah $[t_0, t_0 + T]$ oleh $t_k = t_0 + kT/n$ ($0 \leq k
\leq n$) lalu definisikan $\varphi_n$ secara afin sepotong-sepotong:
$\varphi_n(t_0) = x_0$ dan, pada $[t_k, t_{k+1}]$,
\[
\varphi_n(t) = \varphi_n(t_k) + (t -
t_k)\,f\bigl(t_k, \varphi_n(t_k)\bigr).
\]
\begin{enumerate}
  \item Tunjukkan dengan induksi bahwa $\varphi_n$ terdefinisi dengan baik,
        dengan $\norm{\varphi_n(t) - x_0} \leq M(t - t_0) \leq b$
        pada $[t_0, t_0 + T]$ --- sehingga titik pemasukan nilainya tetap
        di $R$. \emph{(Di sinilah $T \leq b/M$ masuk.)}
  \item Tunjukkan bahwa setiap $\varphi_n$ bersifat Lipschitz-$M$.
  \item Simpulkan dari Arzelà--Ascoli
        (\cref{thm:b3:complete:ascoli}) bahwa ada subbarisan
        $(\varphi_{n_j})$ yang konvergen seragam pada $[t_0, t_0 +
        T]$ ke suatu $\varphi$, yang sendirinya Lipschitz-$M$ dengan
        $\varphi(t_0) = x_0$.
\end{enumerate}

\medskip
\noindent\textbf{Bagian II --- Limitnya menyelesaikan persamaan.}
Definisikan \emph{cacatnya} $\Delta_n(t) = \varphi_n'(t) -
f\bigl(t, \varphi_n(t)\bigr)$ di titik bukan kisi.
\begin{enumerate}[resume]
  \item Tunjukkan bahwa $f$ \omterm{def:b3:topology:continuity}{kontinu} seragam pada $R$, lalu
        simpulkan: untuk setiap $\varepsilon > 0$ ada $n_0$ sedemikian
        sehingga untuk $n \geq n_0$ dan setiap $t$ bukan kisi berlaku
        $\norm{\Delta_n(t)} \leq \varepsilon$. \emph{(Pada $(t_k,
        t_{k+1})$ berlaku $\varphi_n'(t) = f(t_k, \varphi_n(t_k))$,
        sedangkan $(t, \varphi_n(t))$ berjarak paling banyak
        $(1 + M)\,T/n$ dari $(t_k, \varphi_n(t_k))$.)}
  \item Tegakkan bentuk integralnya: untuk setiap $t$,
        \[
        \varphi_n(t) = x_0 + \int_{t_0}^{t}
        f\bigl(s, \varphi_n(s)\bigr)\dd s +
        \int_{t_0}^t \Delta_n(s)\,\dd s,
        \]
        dengan integran di tengahnya \omterm{def:b3:topology:continuity}{kontinu} sepotong-sepotong.
  \item Ambil limitnya sepanjang $(n_j)$: tunjukkan
        $f(s, \varphi_{n_j}(s)) \to f(s, \varphi(s))$
        secara \emph{seragam} (sekali lagi lewat \omterm{def:b3:topology:continuity}{kekontinuan} seragam $f$), lalu
        simpulkan
        \[
        \varphi(t) = x_0 + \int_{t_0}^t
        f\bigl(s, \varphi(s)\bigr)\dd s .
        \]
  \item Simpulkan bahwa $\varphi$ bersifat $\mathcal C^1$ pada $[t_0, t_0 +
        T]$ dan menyelesaikan $x' = f(t, x)$, $x(t_0) = x_0$; lalu perluaslah
        konstruksinya ke $[t_0 - T, t_0]$ (lewat pembalikan waktu).
        \emph{Inilah teorema Peano.}
\end{enumerate}

\medskip
\noindent\textbf{Bagian III --- Ketunggalannya sungguh gagal.}
Tinjaulah $x' = 2\sqrt{\abs x}$, $x(0) = 0$, pada $\R$.
\begin{enumerate}[resume]
  \item Periksalah bahwa $f(x) = 2\sqrt{\abs x}$ \omterm{def:b3:topology:continuity}{kontinu} tetapi
        tidak Lipschitz pada \omterm{def:b3:topology:topology}{persekitaran} $0$ mana pun.
  \item Periksalah bahwa $x \equiv 0$ dan, untuk setiap $c \geq 0$,
        \[
        x_c(t) = \begin{cases} 0 & t \leq c,\\ (t - c)^2 & t
        \geq c,\end{cases}
        \]
        semuanya merupakan penyelesaian yang melalui $(0,0)$: jadi ada kontinum
        penyelesaian yang berbeda.
  \item Di manakah argumen iterasi Picard (lewat titik tetap
        Banach) runtuh untuk $f$ ini?
\end{enumerate}

\medskip
\noindent\textbf{Bagian IV --- Batas metodenya.}
\begin{enumerate}[resume]
  \item Tunjukkan bahwa teorema Peano gagal pada dimensi tak berhingga:
        kita terima (atau Anda boleh percaya saja) contoh klasik
        Dieudonné di ruang $c_0$ berisi barisan yang menuju nol;
        sebagai gantinya, buktikan bahan berdimensi berhingga yang gagal di sana:
        bahwa bola satuan tertutup $c_0$ (dengan norma sup) tidak \omterm{def:b3:topology:compact}{kompak} --- tampilkan
        barisan terbatas yang tak punya subbarisan konvergen, lalu jelaskan
        langkah Bagian I mana yang runtuh.
  \item Rangkumlah: hipotesis mana yang memberikan keberadaan? dan mana yang memberikan keberadaan
        sekaligus ketunggalan? Nyatakan kedua teorema itu secara tepat
        (Peano; Cauchy--Lipschitz) berdampingan.
\end{enumerate}

\medskip
\noindent\textbf{Bagian V --- Osgood: ketunggalan di luar
Lipschitz.} Misalkan $\omega\colon\intoo0\infty\to\intoo0\infty$
\omterm{def:b3:topology:continuity}{kontinu} dan tak turun, dengan
\[
\int_0^1\frac{\dd r}{\omega(r)} = +\infty
\qquad(\text{divergen di } 0),
\]
lalu andaikan $f$ memenuhi $\norm{f(t, x) - f(t, y)} \leq
\omega\bigl(\norm{x - y}\bigr)$ pada $R$.
\begin{enumerate}[resume]
  \item Periksalah bahwa $\omega(r) = Lr$ memenuhinya (yakni Lipschitz), bahwa
        $\omega(r) = r\log\frac1r$ (yang diperluas secara \omterm{def:b3:topology:continuity}{kontinu} untuk
        $r$ kecil) juga memenuhinya walaupun ia \emph{bukan} $O(r)$,
        dan bahwa $\omega(r) = 2\sqrt r$ tidak. Hitunglah
        integralnya pada setiap kasus.
  \item Misalkan $x_1, x_2$ menyelesaikan persamaan itu pada $[t_0, t_0 + T]$
        dengan $x_1(t_0) = x_2(t_0)$, dan $\delta(t) =
        \norm{x_1(t) - x_2(t)}$. Tunjukkan, hanya dari bentuk integralnya,
        bahwa untuk $t_0 \leq s \leq t$ berlaku
        \[
        \delta(t) \leq \delta(s) +
        \int_s^t\omega\bigl(\delta(v)\bigr)\dd v .
        \]
  \item (Teorema Osgood) Andaikan $\delta(t_1) > 0$ untuk suatu
        $t_1$, lalu misalkan $\tau = \sup\{t \leq t_1 : \delta(t) =
        0\}$. Untuk $s \in \intoo\tau{t_1}$, tetapkan $u(t) = \delta(s)
        + \int_s^t\omega(\delta(v))\,\dd v$. Tunjukkan $\delta \leq
        u$, $u' = \omega(\delta) \leq \omega(u)$, lalu simpulkan
        \[
        \int_{\delta(s)}^{u(t_1)}\frac{\dd r}{\omega(r)}
        \leq t_1 - s \leq t_1 - \tau .
        \]
        Biarkan $s \downarrow \tau$ lalu turunkan kontradiksi dengan
        divergennya integral itu. Simpulkan:
        \emph{penyelesaian yang melalui syarat awal yang sama pastilah
        berimpit} --- yakni ketunggalan di bawah syarat Osgood.
  \item Tariklah akibatnya: ketunggalan Cauchy--Lipschitz merupakan
        kasus $\omega(r) = Lr$; persamaan $x' =
        x\log\frac1{\abs x}$ (yang diperluas dengan $0$ di $0$) berpenyelesaian
        tunggal walaupun ruas kanannya tidak Lipschitz di
        $0$; dan untuk $x' = 2\sqrt{\abs x}$ konvergennya
        $\int_0\frac{\dd r}{2\sqrt r}$ persis merupakan hal yang membiarkan sebuah
        penyelesaian meninggalkan $0$ dalam waktu berhingga --- cocokkan nilai
        integral $\int_0^{h^2}\frac{\dd r}{2\sqrt r} = h$
        dengan perilaku lolosnya $x_c$.
\end{enumerate}

\medskip
\noindent\textbf{Bagian VI --- Laju, skema, corong.}
\begin{enumerate}[resume]
  \item (Lema Grönwall integral) Misalkan $e, \eta \geq 0$
        \omterm{def:b3:topology:continuity}{kontinu} pada $[t_0, t_0+T]$ dengan $e(t) \leq \int_{t_0}
        ^t\bigl(L\,e(s) + \eta(s)\bigr)\dd s$ untuk setiap $t$. Tunjukkan
        \[
        e(t) \leq \frac{\sup\eta}{L}\bigl(\eu^{L(t - t_0)} -
        1\bigr)
        \]
        \emph{(tetapkan $G(t) = \int_{t_0}^t(Le + \eta)$, catat bahwa $G'
        \leq LG + \sup\eta$, lalu turunkan
        $\eu^{-Lt}G(t)$)}.
  \item (Euler konvergen dengan laju) Andaikan kini $f$ bersifat
        Lipschitz-$L$ terhadap $x$ dan Lipschitz-$L'$ terhadap $t$ pada $R$.
        Dengan menggabungkan batas cacat pada pertanyaan 4 (yang dibuat
        kuantitatif: $\norm{\Delta_n} \leq (L' + LM)T/n$) dengan
        pertanyaan 17, buktikan
        \[
        \sup_{[t_0, t_0+T]}\norm{\varphi_n - \varphi}
        \;\leq\; \frac{(L' + LM)\,T}{n}\cdot
        \frac{\eu^{LT} - 1}{L} = O\Bigl(\frac1n\Bigr),
        \]
        dengan $\varphi$ berupa \emph{satu-satunya} penyelesaiannya: jadi dengan
        data Lipschitz seluruh barisannya konvergen, dengan laju yang
        eksplisit --- tanpa perlu subbarisan. Mengapa ketunggalannya
        meningkatkan kekonvergenan subbarisan menjadi kekonvergenan penuh
        bahkan tanpa perhitungan ini?
  \item (Skemanya memilih) Untuk $x' = 2\sqrt{\abs x}$, $x(0) =
        0$: tunjukkan bahwa setiap poligon Euler identik nol,
        sehingga skemanya konvergen ke penyelesaian $x \equiv 0$;
        tetapi bila dimulai di $x(0) = \varepsilon > 0$ ia konvergen
        (ketika $n \to \infty$, lalu $\varepsilon \to 0$) ke $t
        \mapsto t^2$, yaitu penyelesaian yang \emph{berbeda} melalui titik
        asalnya. Jadi ketaktunggalannya muncul kembali sebagai kepekaan
        skema numeriknya terhadap usikan.
  \item Tunjukkan bahwa himpunan $\mathcal S$ berisi \emph{semua} penyelesaian
        $x' = f(t,x)$, $x(t_0) = x_0$ pada $[t_0, t_0+T]$
        (dengan nilai di $\bar B(x_0, b)$) bersifat tak kosong (Bagian II),
        Lipschitz-$M$ secara seragam, dan tertutup di
        $\bigl(\mathcal C([t_0, t_0+T], \R^d),
        \norm\cdot_\infty\bigr)$; lalu simpulkan lewat Ascoli bahwa
        $\mathcal S$ bersifat \emph{\omterm{def:b3:topology:compact}{kompak}}. (Teorema Kneser menambahkan
        bahwa $\mathcal S$ \omterm{def:b3:topology:connected}{terhubung}; kita tak akan
        membuktikannya.)
  \item Periksalah gejala Kneser pada contohnya: untuk $x' =
        2\sqrt{\abs x}$, $x(0) = 0$, pada $[0, 1]$, tunjukkan
        $\mathcal S = \{x_c : c \in \intcc0\infty\}$ dengan
        $x_\infty \equiv 0$ \emph{(untuk sembarang penyelesaian, misalkan $c =
        \sup\{t : x(t) = 0\}$ lalu integralkan $(\sqrt x)' = 1$ pada
        $\{x > 0\}$)}, bahwa $c \mapsto x_c$ \omterm{def:b3:topology:continuity}{kontinu} dari
        $\intcc0\infty$ (yakni \omterm{def:b3:topology:locallycompact}{kompaktifikasi satu titiknya}, dengan
        $c = \infty$ direkatkan sebagai limitnya) ke $\mathcal C([0,1])$,
        lalu simpulkan bahwa $\mathcal S$ memang \omterm{def:b3:topology:compact}{kompak} dan
        \omterm{def:b3:topology:connected}{terhubung} --- yaitu corong berbentuk ruas garis.
  \item (Himpunan tercapai) Simpulkan dari pertanyaan 20 bahwa untuk setiap
        $t$ yang tetap, \emph{himpunan tercapai} $\mathcal S(t) =
        \{x(t) : x \in \mathcal S\}$ bersifat \omterm{def:b3:topology:compact}{kompak}; lalu hitunglah untuk
        contoh pada pertanyaan 21 dan periksalah bahwa ia juga
        \omterm{def:b3:topology:connected}{terhubung}: $\mathcal S(t) = \intcc0{t^2}$ --- jadi setiap
        keadaan antaranya dicapai oleh suatu penyelesaian.
\end{enumerate}

\medskip
\noindent\textbf{Bagian VII --- Pelengkap: kebergantungan,
keoptimalan, dan sebuah skema yang dihitung dengan tangan.}
\begin{enumerate}[resume]
  \item (Kebergantungan \omterm{def:b3:topology:continuity}{kontinu}) Andaikan $f$ bersifat Lipschitz-$L$
        terhadap $x$ pada $R$, lalu misalkan $x, y$ dua penyelesaian dengan
        nilai awal $x_0, y_0$ di $t_0$. Dengan menyesuaikan bukti
        pertanyaan 17 pada ketaksamaan $e(t) \leq \norm{x_0
        - y_0} + \int_{t_0}^tL\,e(s)\dd s$, tunjukkan
        \[
        \norm{x(t) - y(t)} \leq \norm{x_0 - y_0}\,
        \eu^{L(t - t_0)},
        \]
        lalu periksalah pada $x' = Lx$ bahwa batasnya tercapai: jadi
        Grönwall bersifat tajam. Simpulkan kembali ketunggalannya (untuk $x_0 =
        y_0$), dan bahwa pemetaan aliran $x_0 \mapsto x(t; x_0)$
        bersifat Lipschitz dengan konstanta $\eu^{LT}$, di mana pun ia
        terdefinisi.
  \item (Osgood bersifat optimal) Sebaliknya, misalkan $\omega$
        \omterm{def:b3:topology:continuity}{kontinu}, tak turun, dan positif pada
        $\intoo0\infty$, dengan $\omega(r) \to 0$ ketika $r \to 0^+$
        tetapi
        \[
        \int_0^1\frac{\dd r}{\omega(r)} < +\infty,
        \]
        lalu perluas $f(x) = \omega(\abs x)$ dengan $f(0) = 0$.
        Tunjukkan bahwa $\Omega(x) = \int_0^x\frac{\dd r}{\omega(r)}$
        merupakan bijeksi naik dari $\intoc0{1}$ pada
        $\intoc0{\Omega(1)}$, bahwa inversnya $g$ menyelesaikan $g'
        = \omega(g)$ dengan $g(0^+) = 0$, dan bahwa $g$, setelah diperluas
        dengan $0$ untuk $t \leq 0$, merupakan penyelesaian $\mathcal C^1$ bagi
        $x' = f(x)$ yang melalui $(0, 0)$ dan berbeda dari $x \equiv
        0$ \emph{(untuk $g'(0) = 0$, batasilah $\frac{g(t)}t$ oleh
        $\omega(g(t))$)}. Simpulkan: hipotesis divergensi
        pada pertanyaan 15 bukanlah kemudahan melainkan perbatasan
        ketunggalan yang persis; lalu perolehlah kembali Bagian III dari
        $\omega(r) = 2\sqrt r$, $\Omega(x) = \sqrt x$.
  \item (Euler dihitung dengan tangan) Untuk $x' = x$, $x(0) = 1$ pada
        $[0, 1]$: tunjukkan bahwa poligon Eulernya memenuhi
        $\varphi_n(1) = \bigl(1 + \frac1n\bigr)^n$. Buktikan
        penjabaran
        \[
        \Bigl(1 + \frac1n\Bigr)^{n}
        = \eu\Bigl(1 - \frac1{2n} + O\Bigl(\frac1{n^2}\Bigr)
        \Bigr),
        \]
        sehingga galatnya di $t = 1$ adalah $\frac{\eu}{2n} +
        O(n^{-2})$: yaitu laju $O(\frac1n)$ pada pertanyaan 18, dengan
        konstanta yang eksak. Periksalah secara numerik untuk $n = 10$:
        $1.1^{10} = 2.59374$ berbanding $\eu \approx 2.71828$,
        yakni galat $0.12454$ yang dibandingkan dengan $\frac{\eu}{20}
        \approx 0.13591$.

\end{enumerate}
\end{problem}
